\documentclass[../main.tex]{subfiles}

\begin{document}

\chapter{Hilbert 空间上的连续线性算子}\label{chapter:ch16-hilbert-operators}

欧氏空间 $\mathbf{R}^{n}$ 中的两个向量 $u=(u_{1},\cdots,u_{n})$ 与 $v=(v_{1},\cdots,v_{n})$ 的内积 $\langle u,v\rangle$ 定义为
\begin{align*}
\langle u,v\rangle=\sum_{k=1}^{n}u_{k}v_{k}.
\end{align*}
我们称此为欧氏内积. 欧氏范数 $\|\cdot\|$ 由关系式
\begin{align*}
\text{对所有 }u\in\mathbf{R}^{n},\quad\|u\|=\sqrt{\langle u,u\rangle}
\end{align*}
确定. 关于欧氏内积有正交性这一重要概念, 它将几何观点带入有限维空间中问题的研究: 子空间有正交补且方程组的可解性由正交关系确定. 内积也带来了一些有十分特殊的结构的有趣的线性算子类: 其中典型的是对称算子, 它有优美的特征向量表示. 本章我们研究具有内积的 Banach 空间 $H$, 如同欧氏空间一样, 该内积与范数相关. 这些空间被称为 Hilbert 空间. 我们证明: 若 $V$ 是 Hilbert 空间 $H$ 的闭子空间, 则 $H$ 是 $V$ 与它的正交补的直和. 基于这种结构性质, 我们证明了 Riesz-Fr\'echet 表示定理, 它刻画了 Hilbert 空间的对偶空间. 由此我们用 Helly 定理推出, Hilbert 空间中的每个有界序列有弱收敛的子序列. 我们证明了 Bessel 不等式, 由此推出可数规范正交集是规范正交基当且仅当它的线性扩张是稠密的. 本章以 Hilbert 空间上的有界对称算子与紧算子的讨论结束, 我们准备证明关于紧对称算子的特征值展开的 Hilbert-Schmidt 定理与关于恒等算子的紧扰动的 Fredholm 性质的 Riesz-Schauder 定理.

\subfile{Chapter-16/section-01.tex}

\subfile{Chapter-16/section-02.tex}

\subfile{Chapter-16/section-03.tex}

\subfile{Chapter-16/section-04.tex}

\subfile{Chapter-16/section-05.tex}

\subfile{Chapter-16/section-06.tex}

\subfile{Chapter-16/section-07.tex}

\end{document}